% =====================================================================
\section{Additional Objectives: Emotion Detection and Language Identification}
\label{sec:extensions}
% =====================================================================

The proposal includes emotion detection (objective~3) and spoken-language
identification (objective~4) \spec. Neither can be trained on GRID alone:
GRID has no emotion labels, and all of its sentences are in English with one
fixed grammar~\cite{cooke2006}. Both are therefore planned as \emph{separate,
optional modules} with their own data, models and evaluation. They share
preprocessing code with the core system but do not change the speech-recognition
architecture or its results, and work on them begins only after the core
milestones in \cref{sec:schedule} are met.

\subsection{Emotion detection from facial movement}
\paragraph{Approach.} The literature reviewed in the proposal selects Face Mesh
landmarks that correspond to Facial Action Coding System regions and converts
them into compact geometric feature vectors, with reported classification
accuracies of up to 97\% on their data~\cite{siam2022}. LipScribe would follow
the same route \tbd:
\begin{enumerate}
  \item reuse the Face Mesh landmarks already computed by the visual pipeline;
  \item compute geometric features per frame---normalised distances (mouth
        width and opening, eyebrow-to-eye distance, eye aperture) and angles
        between landmark triples---after aligning faces to remove head pose and
        scale;
  \item summarise their temporal changes over an utterance with statistics
        (mean, range, velocity) or a small LSTM;
  \item classify with a lightweight model (for example a random forest or
        multilayer perceptron from scikit-learn) into a small set of emotion
        categories.
\end{enumerate}

\paragraph{Data required.} An emotion-labelled audio-visual dataset is
required. Candidates include CREMA-D~\cite{cao2014} and RAVDESS~\cite{livingstone2018},
both of which provide acted emotional speech with frontal video; the team would
confirm licence terms before use. Evaluation would again use speaker-disjoint
splits and report per-class accuracy and a confusion matrix.

\paragraph{Limitations.} Facial movement while speaking is dominated by
articulation, so mouth landmarks partly reflect \emph{what} is said rather than
how the speaker feels. Acted emotions in such datasets are more exaggerated than
natural ones, facial expression is not a reliable indicator of internal emotional
state, and expressions vary across individuals and cultures. Any emotion output
would be presented as an experimental label, not as a statement about the user.

\paragraph{Integration, if completed.} A separate classification head operating
on landmark features, displayed as a tag next to each final caption. It would not
share loss or weights with the speech recogniser.

\subsection{Spoken-language identification}
\paragraph{Approach.} Language identification predicts which language an
utterance is spoken in. A natural design would add a separate utterance-level
classification head to the existing encoders: time-pool the fused representation
$\mathbf{F}$ and apply a linear classifier over the candidate languages, trained
with cross-entropy alongside or after CTC training \tbd. Audio carries most of the
evidence; the proposal's interest is whether lip movements add to it, since visual
speech patterns differ between languages~\cite{aripin2024,kanavnet2026}.

\paragraph{Data required.} Multilingual audio-visual data with language labels is
essential, and recognising speech (not just identifying the language) in each
additional language would need transcribed data and an extended character
vocabulary. Candidate sources are the multilingual MuAViC corpus~\cite{anwar2023}
or language-specific corpora such as those used for Indonesian and Kannada
lipreading~\cite{aripin2024,kanavnet2026}, subject to availability and licensing.
A locally recorded set of short phrases in English and a Ugandan language could
serve as a small qualitative test, with participants' informed consent.

\paragraph{Limitations.} A model trained on GRID has seen one language, a 51-word
vocabulary and studio conditions. It cannot be expected to transfer to other
languages, accents or scripts without new training data. Language identification
on short utterances is also confounded by speaker identity if each language is
spoken by different people, so speaker-disjoint evaluation is again required.

\subsection{Summary of scope}
\begin{table}[H]
\centering
\caption{Core deliverable versus optional extensions.}
\label{tab:scope}
\small
\begin{tabularx}{\textwidth}{@{}p{3.2cm}XXp{2.2cm}@{}}
\toprule
\textbf{Capability} & \textbf{Data} & \textbf{Model change} & \textbf{Status} \\
\midrule
Audio-visual speech recognition & GRID + MUSAN & Full model (\cref{sec:model}) & Core \\
Emotion detection & Emotion-labelled AV corpus (e.g.\ CREMA-D, RAVDESS) & Separate landmark-feature classifier & Optional \\
Language identification & Multilingual AV corpus with language labels & New utterance-level head; new vocabulary for multilingual ASR & Optional \\
\bottomrule
\end{tabularx}
\end{table}
